% =============================================================================
%  Resolución del problema de N-Reinas mediante el algoritmo de Grover
%  Formato: paper científico a dos columnas — 3 hojas A4
% =============================================================================
\documentclass[9pt, a4paper]{extarticle}
\usepackage{multicol}

% --- Codificación y tipografía ---
\usepackage[utf8]{inputenc}
\usepackage[T1]{fontenc}
\usepackage[spanish,es-noshorthands]{babel}
\usepackage{mathpazo}                 % Palatino
\usepackage[scaled=0.88]{helvet}      % Helvética
\usepackage[scaled=0.82]{beramono}    % Bera Mono
\usepackage{microtype}

% --- Matemáticas ---
\usepackage{amsmath, amssymb, amsthm}
\usepackage{physics}

% --- Figuras y tablas ---
\usepackage{graphicx}
\usepackage{float}
\usepackage{booktabs}
\usepackage{caption}
\usepackage{subcaption}
\captionsetup{font=scriptsize, labelfont=bf, labelsep=period, skip=1pt}

% --- Layout (ajustado para 3 hojas) ---
\usepackage[
    top=1.2cm, bottom=1.5cm,
    left=1.1cm, right=1.1cm,
    headheight=16pt,
    footskip=0.6cm
]{geometry}

% --- Colores ---
\usepackage[dvipsnames,svgnames,x11names]{xcolor}
\definecolor{accent}{HTML}{1A3C6E}
\definecolor{accentlight}{HTML}{2D6AB6}
\definecolor{secgray}{HTML}{444444}
\definecolor{rulegold}{HTML}{B8860B}
\definecolor{boxbg}{HTML}{F0F4FA}

% --- Links ---
\usepackage{hyperref}
\hypersetup{
    colorlinks=true,
    linkcolor=accentlight,
    citecolor=accent,
    urlcolor=accentlight,
    pdftitle={N-Reinas via Grover},
    pdfauthor={UNLP}
}

% --- Secciones personalizadas ---
\usepackage{titlesec}
\titleformat{\section}
    {\normalsize\bfseries\sffamily\color{accent}}
    {\thesection.}{0.4em}{\MakeUppercase}
\titleformat{\subsection}
    {\small\bfseries\sffamily\color{secgray}}
    {\thesubsection}{0.4em}{}
\titleformat{\subsubsection}
    {\footnotesize\bfseries\itshape\color{secgray}}
    {\thesubsubsection}{0.4em}{}
\titlespacing*{\section}{0pt}{6pt plus 2pt minus 2pt}{2pt plus 1pt}
\titlespacing*{\subsection}{0pt}{4pt plus 1pt minus 1pt}{1pt}
\titlespacing*{\subsubsection}{0pt}{3pt plus 1pt minus 1pt}{1pt}

% --- Header/Footer ---
\usepackage{fancyhdr}
\pagestyle{fancy}
\fancyhf{}
\renewcommand{\headrulewidth}{0.6pt}
\renewcommand{\headrule}{\hbox to\headwidth{%
    \color{accent}\leaders\hrule height \headrulewidth\hfill}}
\fancyhead[L]{\scriptsize\sffamily\color{secgray}N-Reinas v\'ia Algoritmo de Grover}
\fancyhead[R]{\scriptsize\sffamily\color{secgray}(E1307) Intro.\ Arq.\ Comp.\ Cu\'anticas --- UNLP}
\fancyfoot[C]{\scriptsize\sffamily\color{secgray}\thepage}
\fancypagestyle{plain}{%
  \fancyhf{}%
  \renewcommand{\headrulewidth}{0.6pt}%
  \renewcommand{\headrule}{\hbox to\headwidth{%
    \color{accent}\leaders\hrule height \headrulewidth\hfill}}%
  \fancyhead[L]{\scriptsize\sffamily\color{secgray}N-Reinas v\'ia Algoritmo de Grover}%
  \fancyhead[R]{\scriptsize\sffamily\color{secgray}(E1307) Intro.\ Arq.\ Comp.\ Cu\'anticas --- UNLP}%
  \fancyfoot[C]{\scriptsize\sffamily\color{secgray}\thepage}%
}

% --- Enumerados ---
\usepackage{enumitem}
\setlist{nosep, leftmargin=1.2em}

% --- Decoración ---
\newcommand{\accentrule}{\par\noindent\textcolor{rulegold}{\rule{\columnwidth}{0.8pt}}\par}
\newcommand{\BigO}{\mathcal{O}}

% --- Espaciado ultra-compacto ---
\setlength{\parskip}{0.2pt plus 0.1pt minus 0.1pt}
\setlength{\abovecaptionskip}{1pt}
\setlength{\belowcaptionskip}{0pt}
\setlength{\textfloatsep}{2pt plus 1pt minus 1pt}
\setlength{\floatsep}{2pt plus 1pt minus 1pt}
\setlength{\intextsep}{2pt plus 1pt minus 1pt}
\setlength{\abovedisplayskip}{2pt plus 1pt minus 1pt}
\setlength{\belowdisplayskip}{2pt plus 1pt minus 1pt}
\setlength{\columnsep}{0.45cm}
\raggedbottom

% =============================================================================
\begin{document}

% ======================== TÍTULO ========================
\begin{center}

\vspace*{-0.4cm}

\includegraphics[width=4cm]{LOGO-Ingenieria-color-01-01-1536x587.png}

\vspace{0.02cm}
{\fontsize{7.5}{9}\selectfont\sffamily\color{secgray}%
UNIVERSIDAD NACIONAL DE LA PLATA\; $\cdot$\; FACULTAD DE INGENIER\'IA}

\vspace{0.01cm}
{\fontsize{7.5}{9}\selectfont\sffamily\color{secgray}%
C\'atedra (E1307) Introducci\'on a la Arquitectura de Computadoras Cu\'anticas}

\vspace{0.04cm}
\textcolor{accent}{\rule{\textwidth}{1.5pt}}
\vspace{0.03cm}

{\fontsize{13}{16}\selectfont\bfseries\sffamily\color{accent}%
Resoluci\'on del Problema de N-Reinas\par
\vspace{2pt}
mediante el Algoritmo de Grover}

\vspace{-0.2cm}
\textcolor{accent}{\rule{\textwidth}{0.5pt}}
\vspace{-0.1cm}

{\small\sffamily Ignacio Andres Schwindt (03229/0) $\cdot$ Viernes 17 de Junio de 2026}

\end{center}

\vspace{0.03cm}

\noindent
\colorbox{boxbg}{%
\begin{minipage}{\dimexpr\textwidth-2\fboxsep\relax}
\vspace{2pt}
{\small\sffamily\bfseries\color{accent} RESUMEN}\par\vspace{1pt}
{\small
Se presenta una implementaci\'on del algoritmo de b\'usqueda cu\'antica de Grover
aplicada al problema combinatorio de las N-reinas. La codificaci\'on asigna
$k = \lceil\log_2 N\rceil$ qubits por reina para representar su columna,
eliminando conflictos de fila por construcci\'on. El or\'aculo detecta conflictos
de columna y diagonal mediante enumeraci\'on expl\'icita de patrones y qubits flag,
aplicando fase $-1$ a las configuraciones v\'alidas a trav\'es del esquema
\emph{compute--phase--uncompute}. Se verifica la correcci\'on completa para
$N = 4$ (14~qubits, 8~iteraciones, $P(\text{\'exito}) \approx 99.6\%$) y se
valida el or\'aculo de forma aislada para $N = 5$ (30~qubits), confirmando que
las 10 soluciones conocidas son correctamente identificadas.
Se discuten las limitaciones de escalabilidad y posibles optimizaciones
mediante sumadores cu\'anticos.
}
\vspace{2pt}
\end{minipage}}

\vspace{0.01cm}
\noindent
\colorbox{boxbg}{%
\begin{minipage}{\dimexpr\textwidth-2\fboxsep\relax}
\vspace{2pt}
{\small\sffamily\bfseries\color{accent} PALABRAS CLAVE}\quad
{\small\itshape
Computaci\'on cu\'antica $\cdot$ Algoritmo de Grover $\cdot$ N-reinas $\cdot$
Or\'aculo cu\'antico $\cdot$ Qiskit $\cdot$ B\'usqueda no estructurada}
\vspace{2pt}
\end{minipage}}

\vspace{0.07cm}

\begin{multicols}{2}
\raggedcolumns

% ======================== 1. INTRODUCCIÓN ========================
\section{Introducci\'on}

El problema de las N-reinas consiste en ubicar $N$ reinas en un tablero
de $N \times N$ de modo que ninguna amenace a otra: sin compartir fila,
columna ni diagonal. Formulado originalmente por Max Bezzel en 1848 y
estudiado extensivamente por Gauss, el problema se ha convertido en un
cl\'asico de la teor\'ia combinatoria y la ciencia de la
computaci\'on~\cite{nqueens}.
Para $N = 4$ existen exactamente 2~soluciones; para $N = 5$, hay 10;
para $N = 8$ (tablero est\'andar), 92; y la cantidad crece r\'apidamente.
La b\'usqueda cl\'asica por fuerza bruta tiene complejidad $\BigO(N!)$
al explorar permutaciones de columnas, y aunque algoritmos de backtracking
mejoran esto en la pr\'actica, la complejidad en el peor caso sigue siendo
exponencial.

El \textbf{algoritmo de Grover}~\cite{grover1996}, publicado en 1996,
ofrece una aceleraci\'on cuadr\'atica para problemas de b\'usqueda no
estructurada: si hay $M$ soluciones en un espacio de tama\~no $\mathcal{N}$,
las encuentra en $\BigO\!\big(\sqrt{\mathcal{N}/M}\big)$ consultas al
or\'aculo, contra $\BigO(\mathcal{N}/M)$ en el caso cl\'asico. Esta
aceleraci\'on es \'optima: Bennett, Bernstein, Brassard y Vazirani
demostraron que ning\'un algoritmo cu\'antico puede hacerlo mejor
para b\'usqueda no estructurada~\cite{bennett1997}.

En este trabajo se presenta una implementaci\'on en Python
utilizando Qiskit~\cite{qiskit}, simulada con Aer, que cubre la
codificaci\'on cu\'antica del tablero, el dise\~no del or\'aculo, el
difusor de Grover, y resultados para $N = 4$ (ejecuci\'on completa) y
$N = 5$ (val. or\'aculo).

Adicionalmente, se discute el impacto del mapeo de las restricciones combinatorias hacia compuertas reversibles. Aunque la ejecuci\'on f\'isica de algoritmos con tal profundidad a\'un est\'a limitada por el ruido inherente a los procesadores NISQ (\emph{Noisy Intermediate-Scale Quantum}), esta exploraci\'on establece un caso de estudio \'util sobre c\'omo estructurar y validar or\'aculos de decisi\'on complejos de forma modular.

% ======================== 2. ALGORITMO DE GROVER ========================
\section{El algoritmo de Grover}

\subsection{Principio y flujo general}

Grover opera sobre $n$ qubits con
$\mathcal{N} = 2^n$ estados. El or\'aculo
$U_f$ aplica $U_f\ket{x} = (-1)^{f(x)}\ket{x}$
(\emph{phase kickback} a las soluciones).
El algoritmo (Fig.~\ref{fig:high_level}) tiene tres etapas:
\textbf{(1)}~Inicializaci\'on en superposici\'on uniforme
$\ket{\psi_0} = H^{\otimes n}\ket{0}^{\otimes n}$.
\textbf{(2)}~Iteraci\'on de Grover (repetida $t^*$ veces):
or\'aculo $U_f$ + difusor $D = 2\ket{\psi_0}\!\bra{\psi_0} - I$.
\textbf{(3)}~Medici\'on en la base computacional.

La medici\'on colapsa el estado del sistema, revelando la soluci\'on con una alta probabilidad si se escogi\'o el n\'umero correcto de iteraciones. Este comportamiento probabil\'istico es propio de los algoritmos cu\'anticos, donde la interferencia constructiva amplifica la amplitud de los estados que cumplen con las condiciones del or\'aculo.

\begin{figure}[H]
    \centering
    \includegraphics[width=0.72\columnwidth]{grover_high_level.png}
    \caption{Flujo del algoritmo de Grover para N-reinas.}
    \label{fig:high_level}
\end{figure}

\subsection{Interpretaci\'on geom\'etrica}

El estado del sistema vive en un plano bidimensional
generado por:
\begin{equation}
    \ket{\alpha} = \frac{1}{\sqrt{\mathcal{N}-M}} \sum_{f(x)=0} \ket{x},\quad
    \ket{\beta} = \frac{1}{\sqrt{M}} \sum_{f(x)=1} \ket{x}
\end{equation}
El estado inicial $\ket{\psi_0}$ forma un \'angulo $\theta/2 = \arcsin\sqrt{M/\mathcal{N}}$
con $\ket{\alpha}$. Cada iteraci\'on rota el estado un \'angulo $\theta$
hacia $\ket{\beta}$.

\subsection{Iteraciones \'optimas}

El n\'umero \'optimo de iteraciones es:
\begin{equation}
    t^* = \left\lfloor \frac{\pi}{2\theta} - \frac{1}{2} \right\rceil,\;\;
    \theta = 2\arcsin\!\sqrt{\frac{M}{\mathcal{N}}}
\end{equation}
Para $N\!=\!4$: $\mathcal{N}\!=\!256$, $M\!=\!2$, $\theta \approx 0.177$~rad,
$t^*\!=\!8$, $P(\text{\'exito})\!\approx\!99.6\%$.
\textbf{Nota:} Si se excede $t^*$, la probabilidad \emph{disminuye}
--- el algoritmo ``pasa de largo'' la soluci\'on.
Este fen\'omeno no tiene an\'alogo cl\'asico.

% ======================== 3. CODIFICACIÓN ========================
\section{Codificaci\'on cu\'antica}

\subsection{Eliminaci\'on de la simetr\'ia de fila}

La primera decisi\'on de dise\~no es fijar impl\'icitamente
cada reina~$i$ en la fila~$i$. Esto elimina conflictos de fila
por construcci\'on y reduce el espacio de b\'usqueda de
$N^{2N}$ (posiciones arbitrarias) a $N^N$ (solo columnas por reina).

\subsection{Registro de estado}

La columna de la reina $i$ se codifica en
$k = \lceil\log_2 N\rceil$ qubits binarios. Los qubits
$q_{ik}, q_{ik+1}, \ldots, q_{ik+k-1}$ representan la columna
en base~2, con $q_{ik}$ como bit menos significativo (LSB).
Para $N\!=\!4$: $k\!=\!2$, registro total de $n = Nk = 8$ qubits
(Fig.~\ref{fig:encoding}).

\begin{figure}[H]
    \centering
    \includegraphics[width=0.72\columnwidth]{nqueens_encoding.png}
    \caption{Codificaci\'on de $[1,3,0,2]$ para $N=4$: 2~qubits por reina.}
    \label{fig:encoding}
\end{figure}

\subsection{Qubits ancilla (flags)}

El or\'aculo requiere qubits auxiliares:
\textbf{Flags de pares}: un qubit por cada par $(i,j)$ con $i<j$,
total $\binom{N}{2}$ flags.
\textbf{Flags de rango} (solo cuando $2^k > N$): un qubit por reina
para detectar columnas fuera de $[0, N)$.
Para $N\!=\!4$, $2^2 = 4 = N$, no se necesitan; para $N\!=\!5$
($2^3 = 8 > 5$), se agregan 5 flags.

\begin{table}[H]
    \centering
    \caption{Recursos del circuito seg\'un $N$.}
    \label{tab:resources}
    \scriptsize
    \begin{tabular}{@{}ccccccc@{}}
        \toprule
        $N$ & $k$ & Estado & F.\,pares & F.\,rango & Total & Sols. \\
        \midrule
        4 & 2 & 8  & 6  & 0 & 14 & 2  \\
        5 & 3 & 15 & 10 & 5 & 30 & 10 \\
        6 & 3 & 18 & 15 & 6 & 39 & 4  \\
        \bottomrule
    \end{tabular}
\end{table}

% ======================== 4. ORÁCULO ========================
\section{Dise\~no del or\'aculo}

El or\'aculo es el componente central del circuito.
Aplica fase $-1$ exclusivamente a los estados que representan
configuraciones v\'alidas, en tres etapas (Fig.~\ref{fig:oracle}):
\emph{compute}, \emph{phase flip} y \emph{uncompute}.

\begin{figure}[H]
    \centering
    \includegraphics[width=0.72\columnwidth]{oracle_structure.png}
    \caption{Estructura: Compute $\to$ Phase Flip $\to$ Uncompute.}
    \label{fig:oracle}
\end{figure}

\subsection{Compute Flags}

Para cada par $(i, j)$, $i < j$, se verifica conflicto
($c_i = c_j$ o $|c_i - c_j| = j - i$)
mediante \textbf{enumeraci\'on expl\'icita de patrones}
(Fig.~\ref{fig:conflict}).

\vspace{0.2cm}
\noindent\textbf{Costo Computacional:} Esta verificaci\'on minimiza el ancho aplicando compuertas $X$ \textit{in-place}, pero su profundidad escala como $\BigO(N^3 \log^2 N)$ debido a las Toffoli generalizadas (MCX) requeridas. Esto explica el enorme costo en $N \ge 5$.

\begin{figure}[H]
    \centering
    \includegraphics[width=0.72\columnwidth]{conflict_detection.png}
    \caption{Detecci\'on de conflicto $(c_i, c_j)$ con \texttt{match\_pattern}.}
    \label{fig:conflict}
\end{figure}

\subsubsection{Funci\'on match\_pattern}

Esta es la pieza clave de la detecci\'on. Dado un registro
de $k$ qubits y un valor objetivo $v$, aplica puertas $X$ a cada
qubit cuyo bit correspondiente en $v$ es 0:
\begin{equation}
    \bigotimes_{b=0}^{k-1} X^{1-v_b} \ket{v} = \ket{1}^{\otimes k}
\end{equation}
Esto transforma la condici\'on ``el registro vale $v$'' en
``todos los qubits valen 1'', que es exactamente lo que detecta
una puerta MCX (Toffoli generalizada). Un MCX posterior voltea el
flag correspondiente. Finalmente, se aplica la misma secuencia de $X$
para restaurar el registro ($X^2 = I$).
Para un estado computacional fijo $\ket{c_i, c_j}$, exactamente
una combinaci\'on de patrones coincide, garantizando que el flag se
voltea 0 \'o 1 veces.

\subsubsection{Detecci\'on de fuera de rango}

Para $N$ que no es potencia de 2 (e.g.\ $N = 5$ con $k = 3$), los
valores de columna $\{N, N{+}1, \ldots, 2^k{-}1\}$ son inv\'alidos.
Un flag adicional por reina se activa si su columna codifica un valor
$\geq N$, usando la misma t\'ecnica de \texttt{match\_pattern}.

\subsection{Phase Flip}

Una vez computados todos los flags, se aplica $-1$ condicionada a
que todos valgan 0 (sin conflictos ni columnas fuera de rango):
\begin{equation}
    X^{\otimes m}\!\cdot\!\big(H\!\cdot\!\text{MCX}\!\cdot\!H\big)\!\cdot\!X^{\otimes m}
\end{equation}
donde $m$ es el n\'umero total de flags. La secuencia $X^{\otimes m}$
convierte $\ket{0}^{\otimes m}$ en $\ket{1}^{\otimes m}$ (que dispara el
MCX). El truco $H \cdot \text{MCX} \cdot H$ implementa un MCZ (fase
controlada):
\begin{equation}
    HXH = Z \implies H \cdot \text{MCX}(\vec{c} \to t) \cdot H = \text{MCZ}(\vec{c}, t)
\end{equation}

\subsection{Uncompute Flags}

Para que el or\'aculo sea unitario y reutilizable entre iteraciones,
las ancillas deben volver a $\ket{0}$. Se deshace la etapa de Compute
aplicando las mismas operaciones en orden inverso. Dado que $X$ y MCX
son autoinversas, esta inversi\'on es directa.

\vspace{0.3cm}
\noindent
\colorbox{boxbg}{%
\begin{minipage}{\dimexpr\columnwidth-2\fboxsep\relax}
\vspace{2pt}
{\small\sffamily\bfseries\color{accent} COMPARACI\'ON DE CODIFICACI\'ON}\par\vspace{1pt}
{\small\itshape
A diferencia de formulaciones basadas en simuladores cu\'anticos
de prop\'osito especial con \'atomos en redes \'opticas~\cite{torggler2019},
y de mapeos espaciales que utilizan
$N^2$ qubits asignando uno a cada casillero~\cite{jha2018, santhosh2023},
el enfoque de codificaci\'on por columna presentado aqu\'i requiere solo
$N \lceil\log_2 N\rceil$ qubits de estado. Aunque esto reduce dr\'asticamente
el espacio de Hilbert necesario y garantiza por construcci\'on que no haya
conflictos de fila, conlleva el costo de un or\'aculo de decisi\'on con
mayor profundidad para resolver el resto de las restricciones combinatorias.
}
\vspace{2pt}
\end{minipage}}
\vspace{0.3cm}

% ======================== 5. DIFUSOR ========================
\section{El difusor de Grover}

Implementa la reflexi\'on $D = 2\ket{\psi_0}\!\bra{\psi_0} - I$
(inversi\'on sobre la media) sobre los $n$ qubits de estado.
Su descomposici\'on en puertas es:
\begin{equation}
    D = H^{\otimes n}\!\cdot\!X^{\otimes n}\!\cdot\!\text{MCZ}
        \!\cdot\!X^{\otimes n}\!\cdot\!H^{\otimes n}
\end{equation}
El circuito (Fig.~\ref{fig:diffuser}) opera: (1)~pasar a la base de
Hadamard con $H^{\otimes n}$; (2)~marcar $\ket{0\ldots 0}$ con
$X^{\otimes n}$ seguido de MCZ; (3)~deshacer $X^{\otimes n}$ y volver
a la base computacional.

\begin{figure}[H]
    \centering
    \includegraphics[width=0.68\columnwidth]{diffuser_circuit.png}
    \caption{Circuito del difusor para 8~qubits de estado.}
    \label{fig:diffuser}
\end{figure}

\noindent\textbf{Efecto sobre las amplitudes.} Sea $\bar{a}$ la amplitud media.
El difusor transforma cada amplitud $a_x \to 2\bar{a} - a_x$:
los estados con amplitud por debajo de la media se aten\'uan,
mientras que las soluciones (con amplitud negativa grande) se
reflejan a valores positivos grandes.

% ======================== 6. RESULTADOS ========================
\section{Resultados}

\subsection{\texorpdfstring{$N=4$}{N=4}: ejecuci\'on completa}

Para $N = 4$, el circuito completo se simul\'o con el backend
\texttt{AerSimulator} (m\'etodo \texttt{statevector}) durante 4096~shots.

\begin{table}[H]
    \centering
    \caption{Par\'ametros del circuito para $N = 4$.}
    \label{tab:params}
    \scriptsize
    \begin{tabular}{@{}ll@{}}
        \toprule
        Par\'ametro & Valor \\
        \midrule
        Qubits de estado & 8 (2~bits/columna) \\
        Flags de pares & 6 \\
        Qubits totales & 14 \\
        Espacio de b\'usqueda & $2^8 = 256$ \\
        Soluciones cl\'asicas & 2 \\
        Iteraciones \'optimas & 8 \\
        Profundidad (transpilada) & $\sim$41\,000 \\
        \bottomrule
    \end{tabular}
\end{table}

Las dos soluciones v\'alidas, $[1,3,0,2]$ y $[2,0,3,1]$,
concentran $\sim\!99.6\%$ de la probabilidad total
($\sim$49.6\% cada una, Fig.~\ref{fig:results}).
La decodificaci\'on del bitstring medido se realiza invirtiendo
el orden de bits (Qiskit imprime MSB primero) y reconstruyendo
la columna de cada reina:
$\text{col}_i = \sum_{b=0}^{k-1} c_{ik+b} \cdot 2^b$.

\begin{figure}[H]
    \centering
    \includegraphics[width=0.78\columnwidth]{measurement_results.png}
    \caption{Probabilidad tras 8 iteraciones ($N\!=\!4$). Las dos soluciones
    dominan con $\sim$49.6\% cada una; ruido $<1\%$.}
    \label{fig:results}
\end{figure}

\subsection{\texorpdfstring{$N=5$}{N=5}: validaci\'on del or\'aculo}

El circuito completo para $N\!=\!5$ requiere 30~qubits y $\sim$44
iteraciones ($\sim\!6\!\times\!10^5$ compuertas), fuera del alcance
pr\'actico de la simulaci\'on cl\'asica.
Sin embargo, se valid\'o el or\'aculo de forma aislada, preparando
estados computacionales conocidos y midiendo los flags:

Las \textbf{10 soluciones} cl\'asicas (obtenidas por fuerza bruta
con \texttt{classical\_solutions}) resultaron en 0~flags activos
en todos los casos.
Los casos negativos testeados incluyeron:
\begin{itemize}
    \item $[0,0,0,0,0]$: todas en misma columna.
    \item $[0,1,2,3,4]$: diagonal principal.
    \item $[0,2,4,1,7]$: columna 7 fuera de rango ($\geq 5$).
    \item $[1,3,5,0,2]$: columna 5 fuera de rango.
\end{itemize}
Todos fueron correctamente rechazados con flags $> 0$.
Esto confirma que el or\'aculo separa perfectamente las
configuraciones v\'alidas del resto del espacio de b\'usqueda.

% ======================== 7. COMPLEJIDAD ========================
\section{Complejidad y limitaciones}

\subsection{Complejidad del or\'aculo}

La enumeraci\'on expl\'icita de patrones produce:
$\binom{N}{2} = \BigO(N^2)$ pares de reinas;
hasta $\BigO(N)$ combinaciones conflictivas por par;
cada verificaci\'on: $\BigO(k)$ puertas $X$ m\'as un MCX de $2k$ controles.
Total: $\BigO(N^3)$ llamadas a MCX y $\BigO(N^3 \log N)$ compuertas elementales.
Una implementaci\'on con \textbf{sumadores cu\'anticos} (\emph{quantum adders})
permitir\'ia comparar columnas aritm\'eticamente, reduciendo los MCX
por par de $\BigO(N)$ a $\BigO(\text{poly}(k))$.

\subsection{Escalabilidad}

\begin{table}[H]
    \centering
    \caption{Viabilidad de simulaci\'on seg\'un $N$.}
    \label{tab:scale}
    \scriptsize
    \begin{tabular}{@{}cccl@{}}
        \toprule
        $N$ & Qubits & Iters. & Simulable \\
        \midrule
        4 & 14 & 8    & S\'i (\texttt{statevector}) \\
        5 & 30 & 44   & Solo or\'aculo (\texttt{MPS}) \\
        6 & 39 & $\sim$200 & No (simulador cl\'asico) \\
        \bottomrule
    \end{tabular}
\end{table}

Para $N \geq 5$, las dos v\'ias hacia la ejecuci\'on end-to-end son:
(1)~hardware cu\'antico real con fidelidad y $T_1/T_2$ suficientes
(aspiracional en 2026); o (2)~reescritura del or\'aculo con sumadores
cu\'anticos para reducir la profundidad.

% ======================== 8. CONCLUSIONES ========================
\section{Conclusiones}

\begin{enumerate}
    \item \textbf{Correcci\'on verificada.}
          Para $N\!=\!4$, el algoritmo completo encuentra las 2~soluciones
          con $P \approx 99.6\%$. Para $N\!=\!5$, el or\'aculo clasifica
          correctamente las 10 soluciones y rechaza todos los casos
          inv\'alidos testeados.

    \item \textbf{Aceleraci\'on cuadr\'atica.}
          El algoritmo consulta el or\'aculo
          $\BigO\!\big(\sqrt{\mathcal{N}/M}\big)$ veces, comparado con
          $\BigO(\mathcal{N}/M)$ en b\'usqueda cl\'asica.
          Para $N\!=\!4$: 8~iteraciones vs.\ $\sim$128 consultas cl\'asicas.

    \item \textbf{Dise\~no modular.}
          La separaci\'on en codificaci\'on, or\'aculo
          (compute--phase--uncompute) y difusor permite extender o
          reemplazar componentes de forma independiente.

    \item \textbf{Limitaci\'on principal.}
          La enumeraci\'on expl\'icita genera circuitos profundos
          ($\sim$41\,000 de profundidad para $N\!=\!4$).
          Sumadores cu\'anticos reducir\'ian significativamente la
          profundidad para $N\!\geq\!5$.

    \item \textbf{Perspectiva.}
          La ejecuci\'on en hardware cu\'antico real requerir\'ia
          tiempos de coherencia ($T_1$, $T_2$) muy superiores a los
          disponibles actualmente, adem\'as de correcci\'on de errores
          cu\'anticos. No obstante, la implementaci\'on sirve como
          demostraci\'on conceptual de la ventaja cuadr\'atica de Grover
          sobre problemas combinatorios.
\end{enumerate}

% ======================== BIBLIOGRAFÍA ========================
\begin{thebibliography}{9}
\small
    \bibitem{grover1996}
    L.\,K.~Grover,
    ``\href{https://doi.org/10.1145/237814.237866}{A fast quantum mechanical algorithm for database search},''
    \textit{Proc.\ 28th ACM STOC}, pp.\,212--219, 1996.

    \bibitem{nielsen2010}
    M.\,A.~Nielsen and I.\,L.~Chuang,
    \href{https://doi.org/10.1017/CBO9780511976667}{\textit{Quantum Computation and Quantum Information}},
    Cambridge Univ.\ Press, 10th Ann.\ Ed., 2010.

    \bibitem{qiskit}
    Qiskit Development Team,
    ``\href{https://qiskit.org}{Qiskit: An open-source framework for quantum computing},''
    2024.

    \bibitem{nqueens}
    J.~Bell and B.~Stevens,
    ``\href{https://doi.org/10.1016/j.disc.2007.12.043}{A survey of known results and research areas for $n$-queens},''
    \textit{Discrete Math.}, vol.\,309, no.\,1, pp.\,1--31, 2009.

    \bibitem{bennett1997}
    C.\,H.~Bennett, E.~Bernstein, G.~Brassard, and U.~Vazirani,
    ``\href{https://doi.org/10.1137/S0097539796300933}{Strengths and weaknesses of quantum computing},''
    \textit{SIAM J.\ Comput.}, vol.\,26, no.\,5, pp.\,1510--1523, 1997.

    \bibitem{jha2018}
    R.~Jha, D.~Das, A.~Dash, S.~Jayaraman, B.\,K.~Behera, and P.\,K.~Panigrahi,
    ``\href{https://arxiv.org/abs/1806.10221}{A novel quantum N-queens solver algorithm and its simulation
    and application to satellite communication using IBM Quantum Experience},''
    arXiv:1806.10221 [quant-ph], 2018.

    \bibitem{santhosh2023}
    S.\,G.\,S.~Santhosh, P.~Joshi, A.~Barui, and P.\,K.~Panigrahi,
    ``\href{https://arxiv.org/abs/2312.16312}{A quantum approach to solve N-queens problem},''
    arXiv:2312.16312 [quant-ph], 2023.

    \bibitem{torggler2019}
    V.~Torggler, P.~Aumann, H.~Ritsch, and W.~Lechner,
    ``\href{https://doi.org/10.22331/q-2019-06-03-149}{A quantum N-queens solver},''
    \textit{Quantum}, vol.\,3, p.\,149, 2019.
    arXiv:1803.00735 [quant-ph].
\end{thebibliography}

\end{multicols}

\end{document}
